\section{Conclusion}
\label{sec:conclusion}

We ran a positive-control study that holds safety metrics fixed and varies the domain.
Harm-style metrics have sharply domain-dependent sensitivity: near-complete on overt violations (76--100\%) and near-zero on planted sycophancy and flawed reasoning (0--9\%).
Their null results in those domains are uninformative.
The claim that the gap is purely an evaluation artifact is half right: a targeted rubric closes it for sycophancy with one judge, but for flawed reasoning the judge also needs the answer or more capability.
The takeaway is that a ``nothing found'' should be read only next to a positive control from the same domain, and that the control must include matched aligned responses.

Several follow-ups would strengthen these conclusions.
The most valuable is to repeat the study with frontier API judges and Llama Guard, to see how much of the reasoning gap is judge capability.
Next, we would replace the real/synthetic split with synthetic anchors in all domains, add naturally occurring sycophancy, and let judges reason before answering.
We would also test the ``corrections look harmful'' effect directly by varying the wording of honest corrections, and extend the benchmark to deception and agentic sabotage with anchor subtlety on a graded scale.
